\chapter*{Implicaciones de política pública}
\addcontentsline{toc}{chapter}{Implicaciones de política pública}

\section*{Uno. La consolidación existe en el flujo y no llega al acervo}

Los requerimientos financieros bajan de 4,3\,\% del PIB estimado para 2025 a 4,1 propuestos, y
el balance primario es superavitario en medio punto. Y el saldo histórico sube de 51,4 a
52,3\,\% del PIB.

No hay contradicción: reducir el déficit de un año no reduce lo acumulado. \textbf{Lo que sí
hay es una promesa repetida.} El paquete anterior proyectó el acervo constante en 51,4\,\%
durante seis años y no lo fue: subió nueve décimas en doce meses. Este paquete vuelve a
proyectar estabilidad, un punto más arriba, con la misma estructura de supuestos.

\section*{Dos. El ajuste que sostiene el escenario está fuera del ejercicio}

Los requerimientos llegan a 3,0\,\% del PIB en 2028 y el costo financiero baja un punto entero
de aquí a 2031, suponiendo que la tasa nominal llegue a 5,5\,\% y se quede ahí. La inversión
física baja de 2,5 a 2,3\,\% del PIB y no vuelve.

\textbf{Dos legislaturas separan a este paquete de la fecha en que su propio escenario se
cumple}, y la estabilidad proyectada del acervo descansa sobre una tasa que baja y una
inversión pública que no se recupera.

\section*{Tres. La meta tributaria descansa en administración, no en política}

La recaudación llega a 15,1\,\% del PIB sin subir tasas, y la mitad de su aumento real viene de
aduanas y de cobranza. La elasticidad implícita es de 1,9 contra el cierre estimado.

\textbf{El paquete no separa actividad, eficiencia y miscelánea.} Sin esa separación no hay
manera de juzgar si la cifra es una proyección o una aspiración, y las participaciones a las
entidades ---que crecen 3,7\,\% real--- dependen enteramente de que se cumpla.

\section*{Cuatro. La demografía está haciendo trabajo que el paquete no reconoce}

Las pensiones caen 0,7\,\% real y 4,8\,\% por persona de 65 años y más. La educación crece
3,4\,\% real y 4,4 por persona en edad escolar. \textbf{Los dos denominadores se mueven en
direcciones opuestas y el paquete solo proyecta uno.}

En pensiones, la brecha entre la población que cobra y el gasto que paga convierte un problema
de sostenibilidad en uno de suficiencia. En educación, un dividendo demográfico se está
recibiendo sin contabilizarse, y no se puede aprovechar lo que no se mide.

\section*{Cinco. El etiquetado transversal no cuenta lo que dice contar}

Los doce anexos suman el 74,3\,\% del gasto programable porque 170 de 288 programas están
etiquetados en más de uno. El anexo de igualdad entre mujeres y hombres es en un 48,1\,\%
gasto pensionario.

\textbf{Ninguno de los anexos declara qué fracción de cada programa le corresponde.} El
presupuesto tiene la información y no da el paso, y el resultado es que la principal
herramienta de declaración de prioridades sociales produce totales que no se pueden sumar,
comparar ni auditar.

\section*{Seis, y es de método. Un año de reclasificación exige un diff}

Con 271 programas que aparecen y 358 que desaparecen, cinco ramos que se extinguen, tres que
nacen y una corporación de seguridad que cambia de ramo, \textbf{cualquier lectura de este
paquete que no corra el diff de claves va a publicar recortes que no existen}.

Los tres más grandes de este año: el Ramo 47 cae 95,4\,\% real y no hay decisión de gasto
detrás; Defensa crece 12,3\,\% nominal y sin la Guardia Nacional caería 3,1; y el fondo de
aportaciones federales entero cambia de unidad responsable.

Ese archivo acompaña a este documento y es tan importante como cualquiera de sus capítulos.
\clearpage
